% 【本文件写什么】一级组件 wateros-utils（§ 级聚合）
%   - 组件职责与在内核中的位置：跨组件通用工具（容器、位图、字符串等）。
%   - 聚合 crate os/components/wateros-utils/src/lib.rs 的 pub mod 树与对外导出面
%   - 子模块划分、与相邻组件的依赖边界（谁依赖谁、走哪条门面）
%   - 根 feature / 本组件 Cargo.toml feature 如何选用各 impl
%   - 1～2 段综述后，由下方 \input 展开子模块；不在此逐条列举 api trait
% 【事实来源】
%   - os/components/wateros-utils/src/lib.rs、Cargo.toml
%   - docs/exports/public-api/wateros-utils.md
%   - docs/exports/features/wateros-utils.md
%   - os/feature-tree.txt 中 wateros-utils 相关段
% 【不写什么】
%   - api-v0 契约细节 → 子目录 api/api-v0.tex
%   - 各 impl 算法/平台细节 → 子目录 impl/*.tex

\section{wateros-utils}

\code{wateros-utils}预留与平台弱耦合的通用小工具聚合位置。\code{\#![no\_std]}，无\code{[dependencies]}，无api/impl子crate分层。当前对外仅导出占位函数\code{add(u64, u64) -> u64}，供crate内单测烟测；根\code{wateros}以\code{utils}别名默认依赖，但主线内核代码几乎未调用其公共API。

\code{src/asm/riscv/print\_rigister.S}为早期UART寄存器打印辅助，未通过\code{global\_asm!}或build脚本编入\code{lib.rs}，不参与构建。后续可在此集中纯函数或数据结构，保持不反向依赖\code{wateros-base}等平台类型以维持依赖DAG清晰。
